\subsection{محدودیت‌ها و تحلیل شکست}
علت ثبت دقت \lr{0.00\%} و وضعیت‌های \lr{SKIP\_ARCHITECTURAL} را می‌توان در ناسازگاری بنیادین الگوریتم با لایه‌های پیش‌آموزش‌دیده جستجو کرد:
\begin{enumerate}
    \item \textbf{ناسازگاری فرمول‌بندی بازگشتی با انکودرهای دوطرفه (\lr{DistilBERT, RoBERTa}):} انکودرهای دوطرفه بر پایه انتشار همزمان اطلاعات در تمام جهات آموزش دیده‌اند. تبدیل مکانیزم توجه آنها به یک فرآیند علّی با زوال زمانی یک‌طرفه، مسیر تبادل اطلاعات به سمت چپ را مسدود نموده و منجر به خطای ساختاری شدید در لایه طبقه‌بند نهایی می‌گردد. از این رو سیستم آزمون، آن را به عنوان ناسازگاری ساختاری انکودر (\lr{SKIP\_BERT\_ARCHITECTURAL} و \lr{SKIP\_ROBERTA\_ARCHITECTURAL}) رد کرده است.
    \item \textbf{فروپاشی مقیاس حالت در دیکودر علّی (\lr{GPT-2}):} در مدل \lr{GPT-2}، مخرج نرمال‌سازی کسر $\tilde{\mathbf{Q}}_t \mathbf{Z}_t + \epsilon$ به دلیل ضریب زوال فزاینده در طول دنباله به سمت مقادیر بسیار کوچک میل کرده و نوسانات شدیدی در مقیاس خروجی ایجاد می‌کند که لایه سر طبقه‌بند را دچار اشباع و فروپاشی پیش‌بینی می‌سازد.
    \item \textbf{عدم پشتیبانی از ساختار نوین در \lr{Qwen} (\lr{SKIP\_NO\_SUPPORT}):} عدم سازگاری مستقیم کرنل زوال خطی با ساختار توجه با پرس‌وجوی گروه‌بندی‌شده (\lr{GQA}) و تعبیه مکانی چندبعدی چرخان در کلاس اولیه ماژول، مانع از اجرای آزمون در این مدل شده است.
\end{enumerate}
